\chapter{总结与讨论}

\section{适用范围与额外开销}

本节从使用流程讨论系统可能减少哪些重复操作、增加哪些开销，不作为人工实验基线。第五章实际比较对象是 Qoder，并未测量纯人工完成任务的时间与成本。

\begin{table}[htbp]
\centering
\caption{不同工作方式与本系统的额外开销}\label{tab:benefit-cost}
\begin{tabularx}{\linewidth}{>{\hsize=0.60000\hsize}L>{\hsize=1.20000\hsize}L>{\hsize=1.20000\hsize}L}
\toprule
\textbf{工作方式} & \textbf{本系统希望减少的重复操作} & \textbf{仍需承担的开销} \\
\midrule
人工整理 & 多库查询、重复下载、复制对列、检查重复与冲突，以及补写来源说明 & 模型调用、文件存储、摘要计算、来源适配及人工复核 \\
依靠对话与临时脚本 & 寻找当前文件、补记处理步骤、追查来源及确认哪些结果可以交付 & 维护表结构与解析器、执行检查、保存过程记录和发布版本 \\
\bottomrule
\end{tabularx}
\end{table}

系统并不消除人工，而是让研究者少做重复整理，将注意力放在来源选择、字段含义、冲突与不确定记录上。能减少多少人工工作，仍需另行测量。

\section{对科研数据适用性的实际意义}

发布包将数据表、来源、检查结果与版本记录一并保存。研究者可以按明确的连接键使用多表数据，也能在发现问题时回查并修正。它减少了交付文件缺少说明、不同成员持有不同版本的风险，但是否适合某项分析，仍需结合研究任务判断。

\section{实验结果的关键发现}\label{sec:findings-validation}

十二次列示运行中，十一次正式发布，一次未发布，汇总比例为91.7\%。样例10缺少必需输入，因而没有正式交付。发布状态与需求完成情况必须分别记录；这个比例不代表原始需求全部满足，也不是统一条件下的成功概率。

两案例的 flash 记录成本均是本系统较低，样例6 max 则是 Qoder 较低。交付范围与计量方式不同，无法据此区分执行效率、工作量与其他运行条件的影响。

本系统保存发布清单、来源与检查结果；Qoder 样例6两份离线包也有来源列、清单和方法日志，但缺少所引用的原始文件与脚本。两者便于复核的程度不同，不等于已经测得数值正确率。

\section{核心机制的实际表现}

\textbf{正式发布与文件核验。}成功运行具有可检查文件，未发布运行保留原因。摘要重验确认了既有文件与回执一致，未检验冻结输入后的跨次重跑。

\textbf{运行时表结构检查。}样例7、9展示了动态路线的执行与发布。样例9采用的产品不同于原静态四表规格，不能视作同范围任务已经全部完成；新领域适用范围与变换计算仍需另行检验。

\textbf{证据绑定的人工确认。}样例6 flash 记录三次来源访问授权、两次发布验收及两个版本的替代关系。该记录说明相应步骤实际执行过，但不足以评估模型初审准确率；零 HIL 的动态案例也不证明所有路径均强制触发人工验收。

\section{结论}

\mbox{BioMed QAgent} 将研究需求依次转换为检索、采集、解析、整合与正式交付。Agent 提出查询、来源、多表规格和修改建议，Dataset Core 检查候选并控制正式写入，文件与处理记录随版本保存。

本轮运行展示了可核验交付、缺少必要数据时拒绝发布，以及保留修改历史三项行为。后续应优先核对原始字段与数值、检查明确范围内的参考来源，并在冻结输入、代码和参数后重跑。系统旨在让数据交付有据可查，不以发布成功替代科学判断。

\renewcommand{\bibname}{参考文献}
\begin{thebibliography}{9}
\bibitem{langgraph} LangChain. LangGraph: Interrupts [EB/OL]. \url{https://docs.langchain.com/oss/python/langgraph/interrupts}. 2026-09-06查阅。
\bibitem{autogen} Wu Q., Bansal G., Zhang J., et al. AutoGen: Enabling Next-Gen LLM Applications via Multi-Agent Conversation [EB/OL]. arXiv:2308.08155, 2023.
\bibitem{openai-agents} OpenAI. OpenAI Agents SDK [EB/OL]. \url{https://openai.github.io/openai-agents-python/}. 2026-09-06查阅。
\bibitem{docetl} Shankar S., Parameswaran A. G., Wu E. DocETL: Agentic Query Rewriting and Evaluation for Complex Document Processing [EB/OL]. arXiv:2410.12189v1, 2024.
\bibitem{palimpzest} Liu C., Russo M., Cafarella M., et al. A Declarative System for Optimizing AI Workloads [EB/OL]. arXiv:2405.14696v2, 2024.
\bibitem{temporal-approval} Temporal Technologies. Workflow message passing [EB/OL]. \url{https://docs.temporal.io/encyclopedia/workflow-message-passing}. 2026-09-06查阅。
\bibitem{qoder} Qoder. 什么是Qoder [EB/OL]. \url{https://docs.qoder.com/zh/product-series/what-is-qoder}. 2026-09-06查阅。
\bibitem{biomed-evidence} BioMed-QAgent项目组. Gold6–10 Corrected Campaign Session Report [EB/OL]. 实验代码版本：e5aadfe0；归档目录：\path{docs/evaluation/gold6-10-2026-09-03/}，2026。
\end{thebibliography}

